\documentclass{article}
\usepackage{amsmath,amssymb}
\usepackage{xurl}
\usepackage[hidelinks]{hyperref}
% Shared commands for notes in docs/paper-gaps/.

\DeclareUnicodeCharacter{2080}{\ifmmode{}_{0}\else\textsubscript{0}\fi}
\DeclareUnicodeCharacter{2081}{\ifmmode{}_{1}\else\textsubscript{1}\fi}
\DeclareUnicodeCharacter{2082}{\ifmmode{}_{2}\else\textsubscript{2}\fi}
\DeclareUnicodeCharacter{2099}{\ifmmode{}_{n}\else\textsubscript{n}\fi}

\newcommand{\Lean}{\textsc{Lean}}
\ExplSyntaxOn
\NewDocumentCommand{\leanid}{m}{
  \ifmmode
    \textsf{\detokenize{#1}}
  \else
    \group_begin:
    \urlstyle{sf}
    \tl_set:Nn \l_tmpa_tl {#1}
    \tl_replace_all:Nnn \l_tmpa_tl {\_} {_}
    \exp_args:NV \path \l_tmpa_tl
    \group_end:
  \fi
}
\ExplSyntaxOff
\providecommand{\tr}{\operatorname{tr}}
\newcommand{\tnleanIssueBase}{https://github.com/LionSR/TNLean/issues/}
\newcommand{\tnleanPrBase}{https://github.com/LionSR/TNLean/pull/}
\newcommand{\tnleanDocBase}{https://sirui-lu.com/TNLean/docs/}
\newcommand{\ghissue}[1]{\href{\tnleanIssueBase#1}{GitHub issue \##1}}
\newcommand{\ghpr}[1]{\href{\tnleanPrBase#1}{PR \##1}}
\newcommand{\doclink}[1]{\href{\tnleanDocBase#1.html}{\path{#1}}}

% Dirac notation and the identity operator, as in blueprint/src/macros/common.tex.
% \providecommand keeps note-local definitions authoritative.
\usepackage{dsfont}
\providecommand{\ket}[1]{|#1\rangle}
\providecommand{\bra}[1]{\langle #1|}
\providecommand{\braket}[2]{\langle #1 | #2 \rangle}
\providecommand{\ketbra}[2]{|#1\rangle\!\langle #2|}
\providecommand{\Id}{\mathds{1}}
\providecommand{\one}{\mathds{1}}
\providecommand{\C}{\mathbb{C}}
\providecommand{\MN}[1]{M_{#1}(\C)}
\providecommand{\MD}{M_D(\C)}

% Verdict marker: \gapnote{<kind>}{<status>}. Kind is the policy
% classification (clarification, local-correction, scope-restriction,
% unfaithful, false-source, open-gap); status is open, wip (elimination
% underway), resolved, or historical. Machine-read by the site index;
% prints a verdict line.
\newcommand{\gapnote}[2]{\par\noindent\textbf{Verdict:} #1 (#2).\par}

\interfootnotelinepenalty=10000
\title{Physical-Port Simulation of Bounded Register Channels}
\author{}
\date{2026-10-04}
\begin{document}
\maketitle
\gapnote{scope-restriction}{open}

\section{The source resource convention}
Piroli, Styliaris and Cirac~\cite{Piroli2021LOCC}, Supplement pp.~7--8,
distinguish counted nearest-neighbor gates on fixed-dimensional physical
qudits from free onsite operations involving those qudits and their local
ancillas. In the proof of Proposition~2, the entanglement increase per
intersite layer depends on the physical dimension, independently of the
number of onsite ancillas. This prevents replacing a physical gate by an
arbitrary operation on two enlarged sites at the same unit cost without a
simulation argument.

The existing enlarged-site local-channel model allows arbitrary
intermediate local dimensions. Its composition and trace-norm estimates
remain statements about that model. The construction below supplies a
bounded-register communication theorem; it is not an inclusion theorem
for unrestricted intermediate dimensions.

\section{The proved bounded matching theorem}
Fix a physical dimension $d\ge2$, an upper bound $B$ on the local input
dimension, and $m\le B$. The density witness below forces $m\ge1$, hence $B\ge1$. Set
\begin{align}
 k&=\lceil\log_d B\rceil.
\end{align}
The local basis code into $k$ qudits is chosen independently of the chain
length. Its matrix $E$ is isometric. A completely positive trace-preserving
decoder is obtained by completing $Y\mapsto E^\dagger YE$ outside the code:
\begin{align}
 D(Y)&=E^\dagger YE+
 \operatorname{tr}[(I-EE^\dagger)Y(I-EE^\dagger)]\rho,
\end{align}
where $\rho$ is a fixed density matrix. Then $D(EXE^\dagger)=X$ for every
operator $X$, and the identity persists after tensoring with the identity
on any finite reference. The two-site encoding is the tensor product of
the separate onsite encodings.\footnote{Formal declarations:
\leanid{QuantumCircuit.registerDecoding_comp_encoding},
\leanid{QuantumCircuit.registerDecoding_encoding_reference}, and
\leanid{QuantumCircuit.registerEncoding_pair} in
\doclink{TNLean/Circuit/Channel/RegisterEncoding}.}

On a ring of $N\ge2$ sites, give each site one physical port, a $k$-qudit
data register and a $k$-qudit scratch register. For an arbitrary matching
of CPTP pair channels, extend each channel to its full two-register space
by encoding its output after applying it to a completed decoding of the
input. These extensions are CPTP on the entire register algebra and agree
exactly with the original channels on encoded inputs. There are normalized
finite Kraus representations whose actual placed matching channel has
physical depth at most
\begin{align}
 2k&=2\lceil\log_d B\rceil,
\end{align}
independently of $N$ and the number of selected bonds. Products of $L$
such certified layers have depth at most $2kL$.\footnote{Formal declarations:
\leanid{QuantumCircuit.exists_bounded_matching_simulation},
\leanid{QuantumCircuit.PortRegisters.matchingChannel_isPhysicalPortProtocol},
and \leanid{QuantumCircuit.IsPhysicalPortProtocol.list_prod_uniform} in
\doclink{TNLean/Circuit/Channel/PortSimulation}.}
For $B=1$, the data register has zero qudits and the bound is zero. A
one-site ring's self-loop SWAP is separately proved to be the identity;
the two-distinct-register statement assumes $N\ge2$.

\section{Why the communication cost is uniform}
To exchange one memory digit across a selected bond, swap that digit with
the physical port locally at both endpoints, perform one physical-port
SWAP, and undo the two local swaps. This exchanges the selected memory
digits and restores the ports exactly, including when they carry arbitrary
correlations. Every bond of a matching runs in parallel. Repeating over the
$k$ digits moves each left data register into the right scratch register,
while leaving the right data register in place.

All pair channels can now be executed onsite. Reverse the one common
routing afterward. The total cost is $k+k$, not $2k$ times the number of
bonds. Conjugation by the routing preserves products of channel maps, so
the resulting map is exactly the original placed family on the complete
operator space. Outputs return to their original data registers and all
other wires are restored.\footnote{Formal declarations:
\leanid{QuantumCircuit.IsPhysicalPortUnitary.register_matching} in
\doclink{TNLean/Circuit/Channel/PortRouting} and
\leanid{QuantumCircuit.IsPhysicalPortProtocol.of_routed_family} in
\doclink{TNLean/Circuit/Channel/PortChannelRouting}.}
This construction follows the same gather, execute locally, and return
principle used in the source's proof of Theorem~1 (Supplement pp.~8--9),
but the displayed channel-resource bound is a derived result.

The completed decoder for a joint code need not factor into separated
onsite operations outside that code. No such assertion is used: the joint
extension is applied only after the registers have been colocated. Separate
input/output encoding and decoding use products of local maps. Decoding after encoding
is the identity on every chain operator, with any finite
reference.\footnote{Formal declarations:
\leanid{QuantumCircuit.OnsiteChannel.decodeRegisters_encodeRegisters_reference}
in \doclink{TNLean/Circuit/Channel/OnsiteRegisterEncoding}.}

\section{Remaining comparison and source conditions}
The fixed-port predicate records reduced onsite CPTP operations by local
Kraus families. A full source protocol witness must also account for
locally initialized dilation environments and their disposal, and, when
claiming a pure-state preparation, the source's final system--ancilla
condition. For the asymptotic \textit{QCcc} relation, the source fixes the number
of composed \textit{QCcc} protocols independently of $N$. This is not a
bound on individual local measurement/correction steps: one protocol can
already contain sequential steps at many sites. A source-block decomposition
must also satisfy the within-block control order and one-measurement-per-site
restriction. None of these conditions follows from a quantum-depth bound.

The encoded matching theorem and the product onsite coding identities are
not yet an induction compiling every witness of the existing local-channel
protocol predicate. That comparison must record a common bound on all
intermediate local dimensions, instantiate the global placed-layer
intertwining identities for that predicate, and compose them with site-consistent codes. In
particular, taking $B=B(N)$ after choosing the system size would not supply
the uniform constant claimed above. No equality of phase relations or MPS
classification follows from the matching theorem alone.\footnote{The
resource comparison is tracked by \ghissue{8640}, under \ghissue{8622};
approximate and adaptive conversions are tracked by \ghissue{8380} and
\ghissue{8382}.}

\bibliographystyle{alpha}
\bibliography{references}
\end{document}
